% reunion_6_mar.tex

\begin{frame}
    \begin{itemize}
        \item \fullcite{shirini_multi-objective_2024}
        \item \fullcite{barea_integrated_2024}
    \end{itemize}
\end{frame}

\begin{frame}{\parencite{shirini_multi-objective_2024}}

    \begin{itemize}
        \item Planificación de tiempo y orden de aterrizaje, de un numero dado de aviones en una sola posta de aterrizaje. Utilizando un esquema Multi-Población y Multi-Objetivo (MPMO).
        \item Dos objetivos:
              \begin{itemize}
                  \item Minimizar el consumo de combustible (reduciendo el tiempo de espera).
                  \item Minimizar los costos operativos (penalizaciones por desviaciones del tiempo objetivo de aterrizaje).
              \end{itemize}
        \item \textbf{Funciones objetivo:}
              \begin{itemize}
                  \item \textit{Costo operativo:} Penaliza los aterrizajes tempranos o tardíos (Min).\\ \red{Aterrizar cerca del Target Landing Time (TLT)}
                  \item \textit{Tiempo de espera:} Minimiza el tiempo que los aviones pasan en el aire antes de aterrizar (Min).\\ \red{Aterrizar cerca del EArly Landing Time (ELT)}
              \end{itemize}
    \end{itemize}

\end{frame}
\begin{frame}{\parencite{shirini_multi-objective_2024}}

    \begin{itemize}
        \item No es en tiempo real. Es una planificación estratégica para secuenciar aterrizajes.
        \item No se modela una ruta de aterrizaje como en \parencite{barea_integrated_2024}. Se enfoca en el tiempo de aterrizaje y la secuencia.
        \item \textbf{Método:} Algoritmo genético (NSGA-II) con poblaciones múltiples y estrategias de archivo para mejorar la diversidad de soluciones.
              \begin{itemize}
                  \item MPNSGA-II
                  \item Archive Sharing Strategy (ASS)
              \end{itemize}
    \end{itemize}

\end{frame}


\begin{frame}{\parencite{barea_integrated_2024}}
    \begin{itemize}

        \item \textbf{Problema:} Asignación de pistas y Optimización de trayectorias de descenso.
        \item Es un problema multiobjetivo, pero formulado como uni-objetivo ($\alpha$).
        \item \textbf{Funciones objetivo:}
              \begin{itemize}
                  \item \textit{Retrasos:} Minimiza el tiempo de espera de los aviones.
                  \item \textit{Combustible:} Minimiza el consumo de combustible durante el descenso.
              \end{itemize}
        \item Puede usarse en tiempo real gracias al enfoque de \textit{horizonte rodante}, que divide el problema en intervalos de tiempo más pequeños.
        \item \textbf{Ruta de aterrizaje:} Se modela mediante WPTs. Se optimiza la trayectoria de descenso para cada avión, respetando restricciones de altura y separación.
        \item \textbf{Método:} Modelo matemático no lineal descomponiendo el problema en dos problemas con \textit{descomposición de Benders} y horizonte rodante.
    \end{itemize}



\end{frame}




% \section*{Conceptos Clave}

% \begin{itemize}[leftmargin=*]
%     \item \textbf{Descomposición de Benders:}
%     \begin{itemize}
%         \item Es una técnica para resolver problemas complejos dividiéndolos en partes más pequeñas.
%         \item Se usa cuando el problema tiene dos tipos de variables: unas fáciles de resolver (enteras) y otras difíciles (continuas).
%         \item En este caso, se divide en:
%         \begin{itemize}
%             \item Un problema maestro (asignación de pistas, con variables enteras).
%             \item Subproblemas (optimización de trayectorias, con variables continuas).
%         \end{itemize}
%     \end{itemize}

%     \item \textbf{Horizonte Rodante:}
%     \begin{itemize}
%         \item Es una estrategia para resolver problemas en tiempo real o dinámicos.
%         \item Divide el tiempo en intervalos pequeños (ej. 30 minutos) y optimiza cada intervalo por separado.
%         \item Se usa para adaptarse a cambios en tiempo real, como nuevos aviones que llegan al aeropuerto.
%     \end{itemize}

%     \item \textbf{Ruta de Aterrizaje:}
%     \begin{itemize}
%         \item En el primer paper, no se modela explícitamente. Solo se considera el tiempo de aterrizaje.
%         \item En el segundo paper, sí se modela. Se optimiza la trayectoria de descenso, respetando restricciones de altura y separación entre aviones.
%     \end{itemize}
% \end{itemize}

% \section*{Conclusión}

% Ambos papers abordan problemas clave en la gestión del tráfico aéreo, pero con enfoques diferentes:
% \begin{itemize}
%     \item El primero se centra en la secuencia de aterrizajes y es más estratégico.
%     \item El segundo integra asignación de pistas y trayectorias, con capacidad para operar en tiempo real.
% \end{itemize}

